% Independent derivation. Prefix CFP. Requires amsmath, amssymb, hyperref.
\section{Characteristic-$p$ endomotive maps with the entire split-support lattice}
\label{sec:cc-char-p-full-support}

\paragraph{Original source and conventions.}
The input is Alain Connes, Caterina Consani and Matilde Marcolli,
\href{https://arxiv.org/abs/0806.2401v1}{\emph{Fun with $\mathbb F_1$}},
original \nolinkurl{funBC.tex}, lines 1551--2033, in particular
\texttt{geomfrobprop}/\texttt{sigptaup1},
\texttt{triangular}/\texttt{triangmap}/\texttt{sigptaup}/\texttt{rhoptaup},
\texttt{proptriangular}/\texttt{triangact},
\texttt{filtration}/\texttt{filtr}, and \texttt{triangularbis}.
These source sections were read directly as original LaTeX.
The system originates in Jean-Beno\^it Bost and Alain Connes,
\emph{Hecke algebras, type III factors and phase transitions with
spontaneous symmetry breaking in number theory}, Selecta Math.
(N.S.) 1 (1995), 411--457, cited in that source.
The full carrier is the foundation retained under the byline
The Clankers,
\href{https://github.com/KokunoYumeto/zeta-function-research-reader/blob/a2851f9eb352e7e4376bc629de86fa4ed9c1c434/workbenches/splitzero-tandem/foundations/split-support-geometry-v11/split_support_geometry_arithmetic_curve_v11.tex}{\emph{Split Support Geometry}, v11, Definition \texttt{def:lattice-split}}.
All finite-level maps, support fibers, prime ideals and limiting
maps needed here are proved explicitly below. No purity or
zero-location result is inferred from these algebraic constructions.

Fix a prime $p$, a perfect field $k$ of characteristic $p$, and
an arbitrary nontrivial bounded distributive lattice $L$.
Perfection is used for the invertibility of the coefficient
Frobenius, not for polynomial division or the support arguments.
The complex tropical receiver is replaced only in its ring
coordinate by the actual field $k$; the entire original lattice
and its complete coefficient map remain present.

\subsection{The finite local rings, signs, and the full coefficient map}

For $q=p^\ell$, $\ell\ge0$, use the exact ring
\begin{equation}\tag{CFP1}
 A_q=k[T_q]/(T_q^q-1)=k[\epsilon_q]/(\epsilon_q^q),\qquad
 \epsilon_q=T_q-1,\quad \delta_q=1-T_q=-\epsilon_q.
\end{equation}
The equality follows because $(T-1)^q=T^q-1$ in
characteristic $p$. For $q=1$, this is $k$ and
$\epsilon_1=\delta_1=0$. Every element has its unique
unchanged expression $f=\sum_{j=0}^{q-1}a_j\epsilon_q^j$.
The convolution coordinate in the source is $\delta_q$;
its later filtration coordinate is $\epsilon_q$.
The equality $\delta_q=-\epsilon_q$ will be carried
through each map, including characteristic two where
the two signs coincide in the field.

For any ring $R$ in this section retain
\begin{equation}\tag{CFP2}
\begin{gathered}
 S_R=G_L(R)=\{(a,\lambda):a\in R,\lambda\in L,
                                \ a\ne0\Rightarrow\lambda=1_L\},\\
 (a,\lambda)+(b,\mu)=(a+b,\lambda\vee\mu),\quad
 (a,\lambda)(b,\mu)=(ab,\lambda\wedge\mu),\\
 \tau_R=(0,0_L),\qquad e_R=(0,1_L),\qquad
 \widehat a=(a,1_L),\qquad z_\lambda=(0,\lambda).
\end{gathered}
\end{equation}
For a ring map $f:R\to R'$, its full lift is
$G_L(f)(a,\lambda)=(f(a),\lambda)$. It is well
defined because only a zero amplitude can acquire
a lower label; each operation is preserved directly.
For a $k$-linear map that is not multiplicative the
same formula is additive and $G_L(k)$-linear, with
the identical carrier check; multiplicativity is
never inferred in that case.

Let $K$ be the original tropical semifield with zero,
$K_L=K\otimes_{\mathbb B}L$, and
$\iota_L(\lambda)=1_K\otimes\lambda$.
The nonzero indicator $\beta:K\to\mathbb B$ preserves
sum and product because tropical sums vanish only
when both terms vanish and products vanish only
when one factor vanishes. Thus
$\rho_L=\beta\otimes\operatorname{id}_L$ satisfies
$\rho_L\iota_L=\operatorname{id}_L$. The complete map is
\begin{equation}\tag{CFP3}
 \Phi_R:S_R\hookrightarrow R\times K_L,\qquad
 (a,\lambda)\longmapsto(a,\iota_L(\lambda)).
\end{equation}
It preserves both operations and units. The retraction
recovers the label, so it is injective. Its image
consists exactly of pairs whose second coordinate
lies in $\iota_L(L)$ and equals $1_{K_L}$ whenever
the first is nonzero. For every lifted map $f$,
\begin{equation}\tag{CFP4}
 \Phi_{R'}G_L(f)=(f,\operatorname{id}_{K_L})\Phi_R.
\end{equation}
For linear $f$ this states equality of additive maps;
for ring $f$ it is equality of semiring morphisms.
This distinction preserves the transfer maps below.

The exact top-filtration element is
\begin{equation}\tag{CFP5}
 \widetilde\pi_q=\sum_{j=0}^{q-1}T_q^j
             =\epsilon_q^{q-1}=\delta_q^{q-1},\qquad
 F^iA_q=\epsilon_q^iA_q=\delta_q^iA_q\quad(0\le i\le q).
\end{equation}
To prove the polynomial identity, factor $T^q-1$ in
$k[T]$ as both $(T-1)\sum_{j<q}T^j$ and $(T-1)^q$;
cancellation in the polynomial domain gives the
claimed equality before passing to the quotient.
The sign $(-1)^{q-1}$ is one: $q$ is odd for odd $p$,
and $-1=1$ in characteristic two. In particular no
division by $q$ has been used. For $q>1$,
$\widetilde\pi_q\ne0$ and multiplication by it sends
$\sum a_j\epsilon_q^j$ to $a_0\epsilon_q^{q-1}$.
It detects that coefficient in the top filtration
piece without replacing the filtration by its residue.

\subsection{Transitions and the coefficient Frobenius, with every support fiber}

For $h=p^r$, define the transition and the downward map by
\begin{equation}\tag{CFP6}
\begin{gathered}
 i_{q,hq}:A_q\hookrightarrow A_{hq},\quad
      T_q\longmapsto T_{hq}^h,\quad
      \epsilon_q\longmapsto\epsilon_{hq}^h,\quad
      \delta_q\longmapsto\delta_{hq}^h,\\
 D_{h,q}:A_{hq}\twoheadrightarrow A_q,\quad
      T_{hq}\longmapsto T_q,\quad
      \epsilon_{hq}\longmapsto\epsilon_q,\quad
      \delta_{hq}\longmapsto\delta_q.
\end{gathered}
\end{equation}
The transition sends the basis to the distinct monomials
$\epsilon_{hq}^{hj}$ for $0\le j<q$, proving injectivity.
The downward map is reduction modulo $\epsilon^q$ in
the stated coordinates, and is onto with kernel
$(\epsilon_{hq}^q)$. These facts prove the domains,
codomains and all sign conventions without identifying
the two maps. Transitions compose as
$i_{hq,ghq}i_{q,hq}=i_{q,ghq}$.

On one fixed level define the relative endomorphism
and the coefficient Frobenius by
\begin{equation}\tag{CFP7}
\begin{gathered}
 \sigma_h^{(q)}(\sum a_j\epsilon_q^j)
                     =\sum a_j\epsilon_q^{hj},\qquad
 C_h^{(q)}(\sum a_j\epsilon_q^j)
                     =\sum a_j^h\epsilon_q^j,\\
 \mathfrak F_h^{(q)}=C_h^{(q)}\sigma_h^{(q)},\qquad
               \mathfrak F_h^{(q)}(f)=f^h,\\
 D_{h,q}i_{q,hq}=\sigma_h^{(q)},\qquad
 i_{q,hq}D_{h,q}=\sigma_h^{(hq)}.
\end{gathered}
\end{equation}
The second composite is valid termwise even when the
degree being truncated is at least $q$, because its
image degree at least $hq$ is zero. The first
Frobenius is $k$-linear; $C_h$ acts on coefficients
and is an automorphism because $k$ is perfect.
The commuting maps together give the absolute power
map, exactly the source \texttt{sigptaup1}.
Equivalently on a group monomial with coefficient $a$,
$T_q^j a\mapsto T_q^{hj}a^h=(T_q^ja)^h$.

Put $c=\lceil q/h\rceil\ge1$. Then
\begin{equation}\tag{CFP8}
 \ker\sigma_h^{(q)}=\ker\mathfrak F_h^{(q)}
      =(\epsilon_q^c),\qquad
 \operatorname{im}\sigma_h^{(q)}
      =\operatorname{span}_k\{\epsilon_q^{hj}:hj<q\}.
\end{equation}
Distinct surviving degrees cannot cancel. Coefficient
Frobenius does not change whether a coefficient is
zero, proving both kernels. The image is $k$ when
$h\ge q$ and otherwise is an isomorphic copy of
$A_{q/h}$ by its explicitly listed basis. At $h=1$
the kernel is $(\epsilon^q)=0$ as required.

Every full lift of (CFP6)--(CFP8) keeps the label.
Its power identity is literally
\begin{equation}\tag{CFP9}
 G_L(\mathfrak F_h)(a,\lambda)
       =(a^h,\lambda)=(a,\lambda)^h,
 \qquad
 G_L(C_h)G_L(\sigma_h)=G_L(\mathfrak F_h).
\end{equation}
Finite meets of a repeated lattice label equal that
label; the characteristic-$p$ amplitude power is
additive, and the label of a sum remains its join.
This proves that the power map is the stated actual
semiring morphism, with no disappearing labels.

For any of these maps $G_L(f)$, its three distinct
kernel notions are completely explicit:
\begin{equation}\tag{CFP10}
\begin{gathered}
 G_L(f)^{-1}(\tau)=\{\tau\},\qquad
 G_L(f)^{-1}(e)=\{(a,1_L):a\in\ker f\},\\
 G_L(f)^{-1}(\{z_\lambda:\lambda\in L\})
   =\{z_\lambda:\lambda\in L\}
                  \cup\{(a,1_L):0\ne a\in\ker f\},\\
 G_L(f)(a,\lambda)=G_L(f)(b,\mu)
        \Longleftrightarrow\lambda=\mu\ \text{and }a-b\in\ker f.
\end{gathered}
\end{equation}
The first assertion follows because label zero
forces amplitude zero. The second is the fiber over
supported zero, not an ideal containing $\tau$.
The third is an ordinary ideal and the fourth is
the exact kernel congruence. These statements follow
directly from the two coordinates, and retain their
difference even when a nonzero amplitude is killed.

\subsection{The unnormalized transfer and its exact lifted relations}

There is no inverse of $h=p^r>1$ in $k$. Retain the
source unnormalized transfer, at its exact levels:
\begin{equation}\tag{CFP11}
\begin{gathered}
 R_{h,q}:A_q\longrightarrow A_{hq},\qquad
 R_{h,q}(T_q^j)=\sum_{b=0}^{h-1}T_{hq}^{j+bq},\quad 0\le j<q,\\
 R_{h,q}(\sum_{j<q}a_j\epsilon_q^j)
           =\sum_{j<q}a_j\epsilon_{hq}^{q(h-1)+j},\\
 R_{h,q}(\delta_q^j)=\delta_{hq}^{q(h-1)+j},\qquad
 R_{h,q}(1)=\epsilon_{hq}^{q(h-1)}.
\end{gathered}
\end{equation}
For the second line, the first-line sum is
$T_{hq}^j\sum_{b<h}(T_{hq}^q)^b$; (CFP5) applied
to this polynomial sum gives
$(T_{hq}^q-1)^{h-1}=\epsilon_{hq}^{q(h-1)}$.
Linearity then permits the unchanged polynomial
basis change from $T_q$ to $\epsilon_q$.
For $\delta_q=-\epsilon_q$, the exponent
$q(h-1)$ is even at odd $p$, and its sign is one
at $p=2$, giving exactly the third line.
Distinct resulting degrees prove injectivity.
Its image is the full ideal
$\epsilon_{hq}^{q(h-1)}A_{hq}$, with the displayed
shift providing its inverse as a vector-space map.

For $h>1$ the image is square-zero, because
$2q(h-1)\ge hq$. The map is not a ring map.
The exact projection and composition formulas are
\begin{equation}\tag{CFP12}
\begin{gathered}
 R_{h,q}(D_{h,q}(x)y)=xR_{h,q}(y)
                   \quad(x\in A_{hq},\ y\in A_q),\\
 R_{h,q}D_{h,q}(x)=\epsilon_{hq}^{q(h-1)}x,\qquad
 D_{h,q}R_{h,q}=0\quad(h>1),\\
 R_{g,hq}R_{h,q}=R_{gh,q},\qquad
 i_{hq,bhq}R_{h,q}=R_{h,bq}i_{q,bq}
                                     \quad(b=p^v).
\end{gathered}
\end{equation}
Multiplication by the shift $\epsilon^{q(h-1)}$
kills exactly terms of degree at least $q$ in its
unshifted factor; hence reducing $x$ modulo
$\epsilon^q$ and then multiplying gives the first
two equalities. The downward map kills the image
because $q(h-1)\ge q$. Composition adds the
exponents $hq(g-1)+q(h-1)=q(gh-1)$. The last
identity multiplies every exponent by $b$ on
either side, proving compatibility with all
actual transition levels. All equalities retain
their original coefficients; no factor $1/h$
occurs anywhere.

The lift $\widehat R_{h,q}=G_L(R_{h,q})$ is additive,
injective and $G_L(k)$-linear. Its failure of
multiplicativity and the lifted downward identity
are the following exact equations for $h>1$:
\begin{equation}\tag{CFP13}
\begin{gathered}
 \widehat R_{h,q}(a,\lambda)
           \widehat R_{h,q}(b,\mu)=(0,\lambda\wedge\mu),\\
 G_L(D_{h,q})\widehat R_{h,q}(a,\lambda)=(0,\lambda)
                                      =e_{A_q}(a,\lambda),\\
 \widehat R_{h,q}(G_L(D_{h,q})(x)y)
                   =x\widehat R_{h,q}(y).
\end{gathered}
\end{equation}
Their amplitudes are (CFP12); their labels on
each side are respectively the displayed meet
or unchanged label. Thus even the zero ring
composite is the nonzero supported projection
on the full carrier. The source integral
factor $h$ acts here as the actual scalar
$\widehat h=e$, since $h=0$ in $k$, and
$e(a,\lambda)=(0,\lambda)$ reproduces the
second line. This states what survives that
coefficient, rather than replacing it by $\tau$.

\subsection{Supported nilpotence and the actual full prime spectrum}

For $q>1$ the distinguished generator obeys
\begin{equation}\tag{CFP14}
 \widehat\epsilon_q^{\,q}=e_{A_q}\ne\tau_{A_q},\qquad
 \widehat\epsilon_q^{\,j}\ne e_{A_q}\quad(1\le j<q).
\end{equation}
In fact for every ring $R$ and every positive integer $n$,
$(a,\lambda)^n=(a^n,\lambda)$. Therefore
\begin{equation}\tag{CFP15}
 \{x\in S_R:x^n=\tau_R\text{ for some }n\ge1\}=\{\tau_R\}.
\end{equation}
If the power has label zero, the original label is
zero and its original amplitude is already zero.
No reducedness of $R$ is used. On the other hand,
the elements of $S_{A_q}$ whose powers reach $e$
are exactly the top-labelled elements with amplitude
in $(\epsilon_q)$. For a nonzero such amplitude,
let $d$ be the least index with $a_d\ne0$.
Its $n$th power has leading term
$a_d^n\epsilon_q^{nd}$ when $nd<q$, with nonzero
coefficient because $k$ is a field, and is zero
when $nd\ge q$. Its exact supported nilpotence
index is therefore $\lceil q/d\rceil$.
An amplitude with $a_0\ne0$ is a unit: factor
out $a_0$ and invert its remaining $1+$nilpotent
part by the finite geometric series. Hence
$A_q$ has the unique ring prime $(\epsilon_q)$.

All ordinary semiring prime ideals of $S_{A_q}$ are
\begin{equation}\tag{CFP16}
\begin{gathered}
 \mathfrak M_q=\{z_\lambda:\lambda\in L\}
                  \cup\{(a,1_L):0\ne a\in(\epsilon_q)\},\\
 P_{\mathfrak r}=\{(0,\lambda):\lambda\in\mathfrak r\},
               \qquad\mathfrak r\in\operatorname{SpecLat}L,
 \qquad P_{\mathfrak r}\subsetneq\mathfrak M_q.
\end{gathered}
\end{equation}
To prove exhaustiveness, a prime containing $e$
contains all labelled zeros. Its top amplitudes
form a ring ideal, with additive inverses supplied
by multiplication by $\widehat{-1}$, and primeness
makes it a ring prime. Thus it is $\mathfrak M_q$.
If a prime excludes $e$, it contains no nonzero
amplitude: multiplying such an element by $e$
would put $e$ into it. Its labels form a lattice
ideal prime for meet; absorption gives downward
closure. Conversely the two displayed families
directly satisfy the ideal and prime conditions.
Strictness follows because $e$ is in the first
and in none of the latter. Their lattice
inclusions are precisely those of the original
lattice primes; no Boolean replacement occurs.

The principal supported-zero ideal itself is
$(e)=\{z_\lambda:\lambda\in L\}$. For $q>1$
it is not prime: $\widehat\epsilon_q$ and
$\widehat{\epsilon_q^{q-1}}$ are both outside
it, while their product is $e$. Its radical
is exactly $\mathfrak M_q$, because a power
enters $(e)$ precisely when its amplitude
is nilpotent. For $q=1$ the amplitude ring
is the field $k$, and $(e)=\mathfrak M_1$
is prime. Thus the reduction homeomorphism
below does not preserve the assertion that
the particular principal ideal $(e)$ was
prime: its inverse image is $\mathfrak M_q$,
not merely the original $(e)$.

The topological basic opens are also explicit.
For a nonzero amplitude $a$, $D((a,1_L))$ is
the whole spectrum if $a_0\ne0$, and is exactly
the lattice-prime part if $a_0=0$. For a labelled
zero $z_\lambda$, its basic open is
$\{P_{\mathfrak r}:\lambda\notin\mathfrak r\}$.
These statements follow from (CFP16), including
$D(e)$ equal to the lattice-prime part and
$V(e)=\{\mathfrak M_q\}$.

The intersection of all displayed primes is
$\{\tau\}$. Any nonzero amplitude is excluded
from every lattice prime. For a nonbottom label
$\lambda$, a lattice prime excluding it exists:
choose an ideal maximal subject to being disjoint
from $\uparrow\lambda$; distributivity proves it
prime. Specifically if adjoining either of two
excluded elements meets that filter, while their
meet belongs to the ideal, distributivity puts a
filter element back in the ideal, a contradiction.
Thus $z_\lambda$ is also excluded by a displayed
prime. This verifies the ordinary nilradical
directly from both powers and the actual spectrum.

\subsection{Reduction is a full-label congruence, not a nilradical quotient}

The exact ring reduction and its lift are
\begin{equation}\tag{CFP17}
 \operatorname{red}_q:A_q\to k,\quad
             \sum a_j\epsilon_q^j\longmapsto a_0,
 \qquad
 \mathcal R_q:S_{A_q}\to S_k,\quad
             (a,\lambda)\longmapsto(a_0,\lambda).
\end{equation}
This is onto, with the constant amplitude section
retaining its label. Its kernel congruence is
\begin{equation}\tag{CFP18}
 (a,\lambda)\sim(b,\mu)
      \Longleftrightarrow\lambda=\mu,\ a_0=b_0,
 \qquad
 S_{A_q}/\langle\widehat\epsilon_q=e\rangle\simeq S_k.
\end{equation}
For top labels $a-b=\epsilon_q c$,
multiplying the generating relation by $\widehat c$
and adding $\widehat b$ gives exactly the equality.
Below top label both amplitudes are already zero.
Conversely (CFP17) distinguishes all pairs not
specified, proving the entire congruence.
It is nontrivial for $q>1$ by (CFP14), but its
$\tau$-fiber is only $\{\tau\}$ by (CFP10).
Since the ordinary semiring nilradical is already
$\{\tau\}$, quotienting by that zero ideal leaves
the semiring unchanged and cannot be this map.

Reduction nevertheless induces the exact homeomorphism
\begin{equation}\tag{CFP19}
 \operatorname{Spec}S_k\xrightarrow{\mathcal R_q^*}
          \operatorname{Spec}S_{A_q},\qquad
 \mathfrak M_k\longmapsto\mathfrak M_q,\quad
 P_{\mathfrak r}\longmapsto P_{\mathfrak r}.
\end{equation}
The preimage formulas follow at once from (CFP17).
They give a bijection of the complete list (CFP16).
For each original element, the explicit basic open
calculation following (CFP16) agrees with the
basic open of its reduction: a nilpotent top
amplitude reduces to $e$, not $\tau$. Hence the
bijection and its inverse are continuous.
The unchanged prime topology does not identify
the two coefficient semirings or their congruences.
The faithful complete map instead gives the square
\begin{equation}\tag{CFP20}
\begin{array}{ccc}
 S_{A_q}&\xrightarrow{\Phi_{A_q}}&A_q\times K_L\\
 {\scriptstyle\mathcal R_q}\downarrow&&
      \downarrow{\scriptstyle\operatorname{red}_q\times\operatorname{id}}\\
 S_k&\xrightarrow{\Phi_k}&k\times K_L.
\end{array}
\end{equation}
All labels are retained by both vertical maps;
the nilpotent coordinate changes in the ring
factor exactly as specified. It is a pushout:
on the upper-right semiring the relation
$(\epsilon_q,1)\sim(0,1)$ generates precisely
equal second coordinate and equal reduced ring
coordinate. Multiply it by $(c,0)$ and add
$(b,w)$ to obtain $(b+\epsilon_qc,w)\sim(b,w)$;
reduction proves there are no further relations.
The quotient universal property proves the
stated square, not only its commutativity.

Every transition and relative Frobenius induces
the identity on the label-indexed primes and
takes the unique arithmetic prime to the
unique arithmetic prime: its constant coefficient
is unchanged. Absolute Frobenius raises that
coefficient to a power, which preserves its
vanishing, so it has the same prime action.
These are exact pullback maps on the full
spectrum, although their amplitude kernels
and supported congruences can be different.

\subsection{The convolution limit with the original affine maps}

Let $\Gamma=\mathbb Z[1/p]\cap[0,1)$ and write $b_a$
for the function on $\Gamma$ equal to one at $a$
and zero elsewhere. Define
\begin{equation}\tag{CFP21}
 \mathcal T(p)=\bigoplus_{a\in\Gamma} k b_a,\qquad
 b_ab_c=\begin{cases}b_{a+c}&a+c<1,\\0&a+c\ge1,
             \end{cases}\qquad 1=b_0.
\end{equation}
Only finite supports are allowed. Addition of
the original rational indices and truncation
at 1 prove associativity: a triple survives
exactly when the total of its nonnegative
indices is less than 1, independent of
parenthesization. This is the source
convolution algebra, with $b_a$ denoting its
literal $\delta_a$ to distinguish it from
the finite-ring element $\delta_q$.

The actual compatible maps and their entire
coefficient formulas are
\begin{equation}\tag{CFP22}
\begin{gathered}
 I_q:A_q\hookrightarrow\mathcal T(p),\quad
 T_q\longmapsto b_0-b_{1/q},\quad
 \delta_q\longmapsto b_{1/q},\quad
 \epsilon_q\longmapsto-b_{1/q},\\
 I_q(\sum_{j<q}a_j\epsilon_q^j)
                  =\sum_{j<q}(-1)^j a_j b_{j/q},\qquad
 I_{hq}i_{q,hq}=I_q,\qquad
 \varinjlim_q A_q\xrightarrow{\sim}\mathcal T(p).
\end{gathered}
\end{equation}
For $q=1$, $b_1$ means the truncated zero.
The distinct basis functions prove injectivity.
The sign in the transition agrees because
$(-b_{1/(hq)})^h=-b_{1/q}$ in characteristic $p$.
Every finite set of indices has a common
denominator $q$, proving surjectivity onto
the limit. The maps on $T_q$ therefore prove
the source \texttt{triangmap} with its exact sign.
The source's malformed displayed parenthesis
in \texttt{charpequ1} is read as
$(1-I_p(T_p))^{p-1}=b_{(p-1)/p}$; (CFP5)
and the formula above prove this intended
polynomial identity independently.

The full limit is likewise exact:
\begin{equation}\tag{CFP23}
 \varinjlim_q S_{A_q}\xrightarrow{\sim}S_{\mathcal T(p)},
 \qquad[(a,\lambda)]\longmapsto(I_q(a),\lambda).
\end{equation}
Every amplitude has a representative at a
finite level and the label can be retained
there. Equal amplitudes in the injective
direct limit agree at a common level;
equal elements also have equal labels.
This proves the bijection and its inverse;
finite sum and product use a common level,
so it is a semiring isomorphism. The maps
$\Phi$ commute with it coefficient by coefficient.

For $h=p^r$, the limit Frobenius and transfer
are the following exact affine maps:
\begin{equation}\tag{CFP24}
\begin{gathered}
 \sigma_h(b_a)=\begin{cases}b_{ha}&ha<1,\\0&ha\ge1,
                       \end{cases}\qquad
 (\sigma_h f)(a)=f(a/h),\\
 R_h(b_a)=b_{(a+h-1)/h},\qquad
 (R_h f)(a)=f(ha-(h-1)).
\end{gathered}
\end{equation}
Functions are zero at indices outside $[0,1)$,
as in the source. The transition compatibility
in (CFP12) gives these maps on the limit;
substitution of $a=j/q$ proves their displayed
coefficients directly. In particular the
source formulas \texttt{sigptaup} and
\texttt{rhoptaup} are recovered at $h=p$,
with every original translation $(p-1)/p$.

The exact kernel and image are
\begin{equation}\tag{CFP25}
 \ker\sigma_h=
       \bigoplus_{a\in\Gamma\cap[1/h,1)}k b_a,
 \qquad
 \operatorname{im}R_h=
       \bigoplus_{a\in\Gamma\cap[1-1/h,1)}k b_a.
\end{equation}
For $h=1$ the first is zero and the second
is the whole ring. The surviving map
$a\mapsto ha$ is injective on $[0,1/h)$,
so no cancellation among different surviving
coefficients alters the first formula.
The transfer map is a bijection of the
indicated index sets, proving the second.
The kernel ideal has $h$th ideal power zero,
since $h$ indices at least $1/h$ sum to
at least 1. For $h>1$ the transfer image
has square zero. Its lift has exactly the
support fibers (CFP10) and the product and
composite identities (CFP13), now at the
limit; it is not made multiplicative.

The source filtration elements called
$\tau_m$ in \texttt{basicpropeps} are denoted
$\vartheta_m$ here to keep them separate
from the absorbing split zero $\tau$:
\begin{equation}\tag{CFP26}
 \vartheta_m=R_{p^m}(1)=b_{1-p^{-m}},\qquad m\ge1,
 \quad
 \vartheta_m\vartheta_n=0,
 \quad R_{p^r}(\vartheta_m)=\vartheta_{m+r},
 \quad \sigma_p(\vartheta_m)=0.
\end{equation}
The index sum of two such elements is at
least 1, while the transfer index formula
gives $1-p^{-(m+r)}$ and the Frobenius
index is at least 1. These computations
prove every equality. The positivity of
$m,n$ matters: the additional value
$\vartheta_0=1$ does not obey the asserted
square-zero product law. In the full lift
the first and last zeros in (CFP26) are
$e$, never $\tau$, since all those
displayed elements have top support.

The limit ring is local with maximal ideal
$\mathfrak m=\{f:f(0)=0\}$. If a nonzero
$f\in\mathfrak m$ has least supported
index $a>0$, its powers vanish exactly
starting at $\lceil1/a\rceil$, by their
nonzero leading coefficient before that
index reaches 1. If $f(0)\ne0$, the finite
geometric inverse after factoring that
constant proves it a unit. Thus every
element of $\mathfrak m$ is nilpotent,
but the ideal is not nilpotent: for each
$N$ choose $q>N$, so $b_{1/q}^N\ne0$.
The full prime spectrum and reduction
$S_{\mathcal T(p)}\to S_k$ are exactly
(CFP16)--(CFP20), with $\mathfrak m$
in place of $(\epsilon_q)$. Their proofs
use this actual local-ring description.

\subsection{The prime-to-$p$ reduction with its exact remaining Frobenius}

For $n=qm$, $q=p^\ell$ and $(m,p)=1$, set
\begin{equation}\tag{CFP27}
 R_n=k[T]/(T^n-1)
        =k[T]/((T^m-1)^q),\qquad
 B_m=k[T]/(T^m-1).
\end{equation}
The derivative $mT^{m-1}$ is coprime to
$T^m-1$, so the latter polynomial has
no repeated irreducible factor over $k$.
Its distinct irreducible factors $u_i$
give by polynomial Bezout and the Chinese
remainder maps
\begin{equation}\tag{CFP28}
 R_n\simeq\prod_i k[T]/(u_i^q),\qquad
 B_m\simeq\prod_i k[T]/(u_i).
\end{equation}
Each right factor in the second product
is a field, hence their product is
reduced. The kernel $(T^m-1)$ of
$R_n\to B_m$ has $q$th power zero.
Any nilpotent maps to zero in the
reduced target, proving that this
kernel is exactly the ring nilradical.
The ring primes are the ideals $(u_i)$.
The full semiring primes are those
arithmetic primes together with every
lattice prime $P_{\mathfrak r}$;
each lattice prime is contained in
each arithmetic prime, by the same
two-case proof as (CFP16).

The full reduction keeps labels:
\begin{equation}\tag{CFP29}
 S_{R_n}\longrightarrow S_{B_m},\qquad
 (f,\lambda)\longmapsto([f],\lambda),\qquad
 S_{R_n}/\langle\widehat{T^m-1}=e\rangle\simeq S_{B_m}.
\end{equation}
The equal-label congruence proof of
(CFP18) applies to the displayed
principal ideal. It has trivial
$\tau$-fiber and induces the same
full prime-spectrum homeomorphism,
now with each arithmetic prime
matched to the corresponding field
factor. Ring-prime homeomorphism
can also be checked directly because
every prime contains the nilpotent
kernel; the displayed lattice primes
are unchanged. The complete $\Phi$
square is obtained by keeping $K_L$
fixed exactly as in (CFP20).

Taking the original group inclusions
as transitions gives the source
reduction and its full lift
\begin{equation}\tag{CFP30}
\begin{gathered}
 A=k[\mathbb Q/\mathbb Z],\quad
 B=k[(\mathbb Q/\mathbb Z)^{(p)}],\qquad
 A_{\mathrm{red}}\simeq B,\\
 G_L(A)\longrightarrow G_L(B),\qquad
             (a,\lambda)\longmapsto(\operatorname{red}(a),\lambda).
\end{gathered}
\end{equation}
Indeed every finite-support group-ring
element lies in a finite cyclic level.
The prime-to-$p$ levels $B_m$ are reduced
by (CFP28), and their group-basis
inclusions are injective. Their union
$B$ is therefore reduced. The canonical
primary decomposition of $\mathbb Q/\mathbb Z$
gives $A=\mathcal T(p)\otimes_k B$;
augmenting the first factor to $k$
is the indicated reduction. Its kernel
elements each occur in one finite
level, where they are nilpotent by
(CFP27)--(CFP28). Conversely every
nilpotent maps to zero in $B$. This
proves the full ring reduction and
its exact lifted equal-label congruence.
The ordinary semiring nilradical
remains only $\{\tau\}$ by (CFP15).

On the remaining prime-to-$p$ group,
multiplication by $p$ is an automorphism.
At every finite level $m$ its inverse
is multiplication by the inverse of
$p$ modulo $m$; these inverses agree
under the group inclusions, because
the inverse of a group automorphism
is unique. Thus $\sigma_p$ and all
$\sigma_{p^r}$ become actual
automorphisms on $B$ and on $G_L(B)$,
with labels unchanged. This is the
exact surviving map in source
\texttt{triangularbis} and \texttt{nodenom}.
It is a different transfer after
reduction: the unreduced $R_p$
descends to the zero additive map
on amplitudes, since its image has
positive convolution index, while
the inverse of reduced $\sigma_p$
is unital. Their full lifts are
respectively $(a,\lambda)\mapsto(0,\lambda)$
and the stated label-preserving
automorphism. This proves their
actual relation instead of silently
identifying them. No inverse of $p$
in the field has been introduced.

The finite polynomial quotient and the canonical
prime-to-$p$ group coordinate have the following
exact map, so no generator is silently renamed
in taking that limit. At $n=qm$, choose integers
$a,b$ with $aq+bm=1$. Write $U_m=e(1/m)$ and
$T_n=e(1/n)$ for the actual group generators.
Then
\begin{equation}\tag{CFP31}
 \frac1{qm}=\frac a m+\frac b q\pmod{\mathbb Z},\qquad
 \operatorname{red}(T_n)=U_m^a,\qquad
 k[T]/(T^m-1)\xrightarrow{\sim}k[\mathbb Z/m\mathbb Z],
                         \quad[T]\longmapsto U_m^a.
\end{equation}
The exponent $a$ is the inverse of $q$ modulo $m$;
different Bezout choices give the same power.
It is prime to $m$, proving that the last map
is an isomorphism. This is exactly the primary
group projection, which commutes with every
original group inclusion. On $\mathcal T(p)\otimes B$
the unreduced transfer has the explicit factor map
$R_p\otimes\sigma_p^{-1}$: a preimage in the
prime-to-$p$ group is unique, while the $p$
preimages in the primary group are exactly
those summed by $R_p$. Its first factor has
positive convolution index. Augmentation therefore
makes the whole transfer zero, verifying the
last comparison above on every original basis
element. The reduced inverse remains a different
unital map with the exact factor described there.

\subsection{The original triangular representation and its faithful full lift}

Let $V=k^{(\Gamma)}$ with basis $\xi_a$, $a\in\Gamma$,
and put $\xi_c=0$ for $c\ge1$. Denote convolution by
$f\in\mathcal T(p)$ by $M_f$. The source's two
operators and their exact rational indices are
\begin{equation}\tag{CFP32}
 M_{b_c}\xi_a=\xi_{a+c},\qquad
 U\xi_a=\xi_{(a+p-1)/p},\qquad
 W\xi_a=\xi_{pa}.
\end{equation}
They are respectively the source images of
$\delta_c$, $\widetilde\mu_p$ and $\mu_p^*$.
The star in the last source symbol is its
generator notation; no Hilbert adjoint over
the characteristic-$p$ field is asserted.
Each displayed partial affine map is injective
on the set of indices on which it survives,
so its matrix has at most one nonzero entry
per row and column. It is lower triangular
in the original rational order because
$a+c\ge a$, $pa\ge a$, and
$(a+p-1)/p-a=(p-1)(1-a)/p>0$.
This does not enumerate the dense rational
order increasingly by natural numbers.

The complete defining relations are
\begin{equation}\tag{CFP33}
 UM_fW=M_{R_pf},\qquad
 WM_f=M_{\sigma_pf}W,\qquad
 M_fU=UM_{\sigma_pf},\qquad WU=0.
\end{equation}
For $f=b_c$, the first equality sends $\xi_a$
to $\xi_{a+(c+p-1)/p}$ exactly when
$pa+c<1$, and to zero otherwise, on both
sides. The second sends it to $\xi_{p(a+c)}$
exactly when $p(a+c)<1$. The third sends
it to $\xi_{(a+pc+p-1)/p}$ exactly when
$a+pc<1$. Any earlier cutoff implies the
same failed final inequality, so none is
omitted. Finally $WU\xi_a=\xi_{a+p-1}=0$.
Linearity proves all statements for every
finite-support $f$.

Let $\mathcal C_p$ be the source algebra generated
by $\mathcal T(p)$ and $u,v$ with these relations.
They give its complete normal monomial basis
\begin{equation}\tag{CFP34}
 u^n b_c\quad(n\ge1,c\in\Gamma),\qquad
 b_cv^m\quad(m\ge0,c\in\Gamma).
\end{equation}
Here is a proof of both spanning and independence,
so the use of the source representation is exact.
For $f,g\in\mathcal T(p)$ the products of these
forms reduce by the displayed defining relations to
\begin{equation}\tag{CFP35}
\begin{aligned}
 (u^nf)(u^rg)&=u^{n+r}\sigma_p^r(f)g,\\
 (fv^m)(gv^k)&=f\sigma_p^m(g)v^{m+k},\\
 (fv^m)(u^ng)&=0\quad(m,n\ge1),\\
 f(u^ng)&=u^n\sigma_p^n(f)g,\\
 u^n(fg)v^m&=
 \begin{cases}R_p^n(fg)v^{m-n}&m\ge n,\\
 u^{n-m}R_p^m(fg)&n>m.
 \end{cases}
\end{aligned}
\end{equation}
The first, second and fourth repeatedly commute
the coefficient by the corresponding relation.
The third uses $vu=0$. For the last, apply
$u h v=R_p(h)$ exactly $\min(m,n)$ times.
Expanding the resulting finite-support functions
in their $b_c$ basis proves that the span
contains all generators and is closed under
products, hence spans the generated algebra.

The actions of the normal monomials are
\begin{equation}\tag{CFP36}
\begin{aligned}
 \theta(u^nb_c)\xi_a
    &=\xi_{p^{-n}a+1-p^{-n}+p^{-n}c},\qquad n\ge1,\\
 \theta(b_cv^m)\xi_a&=\xi_{p^m a+c},\qquad m\ge0.
\end{aligned}
\end{equation}
They follow by iterating (CFP32), with the
same cutoff at 1. The first slopes are less
than one, the second at least one. Within
each family its slope and intercept recover
its parameters. Thus distinct normal monomials
have distinct affine functions $t_ja+c_j$,
where $t_j>0$ and $0\le c_j<1$.
For a finite nontrivial linear combination
choose $\eta>0$ smaller than 1 and every
$(1-c_j)/t_j$. The interval $(0,\eta)$
contains infinitely many points of $\Gamma$.
Two distinct affine functions meet at at
most one point, so choose one such point
outside all their finitely many crossings.
Every output then survives and the output
indices are pairwise distinct. A nontrivial
linear combination cannot annihilate that
basis vector. This proves independence and
the faithful unital representation
\begin{equation}\tag{CFP37}
 \theta:\mathcal C_p\hookrightarrow\operatorname{End}_k(V),
 \qquad
 G_L(\theta):G_L(\mathcal C_p)
                 \hookrightarrow G_L(\operatorname{End}_k(V)).
\end{equation}
The lift keeps the original global label.
The formula (CFP2) also defines a semiring
for a noncommutative ring, using its actual
product; the preceding injectivity follows
from the two coordinates. In this full
receiver the source relation is precisely
\begin{equation}\tag{CFP38}
 (v,1_L)(u,1_L)=(0,1_L)
       =\underbrace{(1,1_L)+\cdots+(1,1_L)}_{p\text{ summands}}.
\end{equation}
The top-supported copy of the source ring
has its own additive zero $e$, and its
inclusion in the full carrier does not
send that zero to $\tau$.

\subsection{A coefficientwise full-lattice matrix algebra receiving both actions}

The global receiver is not automatically the
sparse matrix algebra over $G_L(k)$. Construct
an explicit common ambient semiring as follows.
Let $\mathfrak M_L(\Gamma)$ consist of pairs
$(A,\Lambda)$, where $A$ is a column-finite
$k$-matrix on $\Gamma$, $\Lambda:\Gamma^2\to L$
has finite image, and
$A_{ab}\ne0\Rightarrow\Lambda_{ab}=1_L$.
Every entry is thus an actual split coefficient;
the support labels may be dense in the index set.
Define
\begin{equation}\tag{CFP39}
\begin{gathered}
 (A,\Lambda)+(B,\Omega)=(A+B,\Lambda\vee\Omega),\qquad
 (A,\Lambda)(B,\Omega)=(AB,\Lambda\odot\Omega),\\
 (\Lambda\odot\Omega)_{ab}
        =\bigvee_{c\in\Gamma}\Lambda_{ac}\wedge\Omega_{cb},\qquad
 I_{\mathrm{sp}}=(I,\Delta),\quad
 \Delta_{ab}=\begin{cases}1_L&a=b,\\0_L&a\ne b.\end{cases}
\end{gathered}
\end{equation}
The indicated join uses only finite joins in
the original $L$: its values belong to the
finite set of meets of one value from each
input's finite image, and one joins only the
distinct values actually attained. All output
entries lie among finitely many possible
joins of subsets of that set. Hence the
output again has finite image. The amplitude
product is column finite because each column
of $B$ selects only finitely many finite
columns of $A$. Each product entry has a
finite amplitude sum; a nonzero sum has a
nonzero summand and hence a top support path,
proving the carrier condition.

Associativity of the label product follows
by expanding either bracketing as the join
of $\Lambda_{ac}\wedge\Omega_{cd}\wedge\Xi_{db}$.
There are only finitely many distinct such
values. Finite distributivity expands each
side and gives the same set of values; the
same reasoning proves distributivity over
addition. Amplitude operations are those of
column-finite endomorphisms. The zero is
$(0,0)$ and the unit is $I_{\mathrm{sp}}$,
as direct substitution proves. This establishes
the semiring without assuming that arbitrary
joins exist in $L$. Coefficient matrices with
infinitely many distinct labels are not
silently included in this defined ambient
space; it includes every global label and
every sparse original source operator.

Let $\mathbf J$ denote the constant-top matrix
on $\Gamma^2$, and set $H=(I,\mathbf J)$.
The source amplitudes remain the original
triangular matrices in the following exact
corner and faithful coefficient receiver:
\begin{equation}\tag{CFP40}
\begin{gathered}
 H^2=H,\qquad
 H(A,\Lambda)H=(A,\operatorname{const}\epsilon(\Lambda)),
        \quad\epsilon(\Lambda)=\bigvee_{a,b}\Lambda_{ab},\\
 H\mathfrak M_L(\Gamma)H
       \simeq G_L(\operatorname{End}_k(V)),\quad
           (A,\lambda)\longmapsto(A,\operatorname{const}\lambda),\\
 G_L(\mathcal C_p)\hookrightarrow H\mathfrak M_L(\Gamma)H,
           \qquad(x,\lambda)\longmapsto
                         (\theta(x),\operatorname{const}\lambda).
\end{gathered}
\end{equation}
The nonempty index set and finite-image
join formula show that both sides of the
first product have the stated constant
label. If $A\ne0$, one entry forces
$\epsilon(\Lambda)=1_L$. Conversely every
constant-labelled pair obeying that carrier
condition is fixed by $H$ on both sides.
Constant labels multiply by meet and add
by join, giving the exact isomorphism.
Its unit is $H$, which differs from the
ambient sparse unit at every off-diagonal
zero amplitude. Its additive zero is the
ambient $(0,0)$; the image of global $e$
is $(0,\mathbf J)$. No identification of
these three elements has been made.

For an exact action on the original finite
amplitudes, let $\mathcal W_L$ consist of
$(w,\alpha)$ with $w\in V$, $\alpha:\Gamma\to L$
of finite image, and
$w_a\ne0\Rightarrow\alpha_a=1_L$.
Scalar action is $(c,\lambda)(w,\alpha)
=(cw,(\lambda\wedge\alpha_a)_a)$.
The matrix action and its corner are
\begin{equation}\tag{CFP41}
\begin{gathered}
 (A,\Lambda)(w,\alpha)
    =\left(Aw,\left(\bigvee_b\Lambda_{ab}\wedge\alpha_b\right)_a\right),\\
 H(w,\alpha)=(w,\operatorname{const}\bigvee_a\alpha_a),\qquad
 H\mathcal W_L\simeq G_L(V),\\
 (A,\lambda)(w,\mu)=(Aw,\lambda\wedge\mu)
                         \quad\text{on the corner}.
\end{gathered}
\end{equation}
The amplitude remains finitely supported
because $w$ is and $A$ is column finite.
The label joins have finitely many distinct
values, remain in the original lattice,
and satisfy the carrier test by a nonzero
contributing amplitude. Finite distributivity
proves module and action associativity.
Testing the vector with amplitude $\xi_b$
and only label $b$ top extracts the entire
$b$th column, so the ambient action is
faithful. On the global corner, testing
basis amplitudes and the vector $(0,1_L)$
likewise recovers every amplitude and label.

The complete coefficient map also exists
entrywise on this ambient space:
\begin{equation}\tag{CFP42}
 (A,\Lambda)\longmapsto
      (A,(\iota_L(\Lambda_{ab}))_{a,b\in\Gamma}).
\end{equation}
In the support coordinate use finite-image
matrices over $K_L$, with multiplication
defined by the same finite sums of distinct
values. Its idempotent addition permits
the repeated-value interpretation used
above. The morphism preserves every product
and sum because $\iota_L$ preserves their
finite lattice operations. The retraction
$\rho_L$ applied entrywise proves injectivity.
It takes the constant corner label to its
constant $\iota_L$ value, and takes every
support defect below to its exact indicator.
Thus this is the original full $\Phi$ on
every coefficient, with no chosen support
character or lost intermediate label.

\subsection{The exact sparse defects, supported nilpotence, and a connecting morphism}

For an amplitude matrix define its sparse
lift by $\operatorname{sp}(A)=(A,\mathbf1_{\operatorname{supp}A})$,
with actual top/bottom labels. For any
set of entries $D$, write
$Z_D=(0,\mathbf1_D)$. The exact addition
and product comparison is
\begin{equation}\tag{CFP43}
\begin{gathered}
 D_+(A,B)=(\operatorname{supp}A\cup\operatorname{supp}B)
                               \setminus\operatorname{supp}(A+B),\\
 D_\times(A,B)=(\operatorname{supp}A\circ\operatorname{supp}B)
                               \setminus\operatorname{supp}(AB),\\
 \operatorname{sp}(A)+\operatorname{sp}(B)
             =\operatorname{sp}(A+B)+Z_{D_+(A,B)},\\
 \operatorname{sp}(A)\operatorname{sp}(B)
             =\operatorname{sp}(AB)+Z_{D_\times(A,B)}.
\end{gathered}
\end{equation}
Relational composition means the existence
of an intermediate index with both entries
nonzero. The sum's label is top on the
union of supports, and its extra zero
amplitudes are precisely the first defect
set. A product's label is top precisely
where a nonzero path exists, while its
amplitude may cancel after summing paths;
the second defect set is exactly those
remaining entries. This proves every
coefficient of both identities.
In particular $\operatorname{sp}(A)+
\operatorname{sp}(-A)=Z_{\operatorname{supp}A}$
for nonzero $A$, so sparse lifting is
not an additive ring map.

The source gives a sharper example that
retains the characteristic-$p$ factor:
\begin{equation}\tag{CFP44}
\begin{gathered}
 W_{\mathrm{sp}}U_{\mathrm{sp}}=(0,0),\qquad
 pI_{\mathrm{sp}}=(0,\Delta),\\
 (HW_{\mathrm{sp}}H)(HU_{\mathrm{sp}}H)
          =(0,\mathbf J)=pH.
\end{gathered}
\end{equation}
There is no surviving sparse path in
$WU$, by (CFP33), so its support is
bottom everywhere. The $p$ additions
of the sparse identity have zero
amplitude and retain only diagonal
top labels. In the constant corner,
both operators have top label on
all pairs, so their product retains
it even when its amplitude vanishes.
This is exactly the faithful image
of (CFP38), and proves why the two
coefficient presentations cannot be
identified by an unstated lift.

Likewise, for $0<c<1$ in $\Gamma$,
translation $M_{b_c}$ has its first
zero power at $n=\lceil1/c\rceil$.
Before that power the index $0$
survives; at that power every path
ends at an index at least 1.
Consequently
\begin{equation}\tag{CFP45}
\begin{gathered}
 \operatorname{sp}(M_{b_c})^n=(0,0),\qquad
        (M_{b_c},\mathbf J)^n=(0,\mathbf J),
                         \quad n=\lceil1/c\rceil,\\
 \widehat{b_c}^{\,n}=e_{\mathcal T(p)}\ne\tau_{\mathcal T(p)}.
\end{gathered}
\end{equation}
The scalar assertion (CFP15) therefore
does not claim that every sparse
matrix semiring lacks nonzero
nilpotents to its ambient zero.
The sparse matrices have them;
the displayed full global lift
retains their original supported
nilpotence instead. Equations
(CFP40)--(CFP45) are their exact
connecting maps and defects.

Saturation $M\mapsto HMH$ is additive
but is not generally multiplicative.
For sparse matrix units $E_{ij}$ and
$E_{kl}$ with $j\ne k$, their product
has no path and saturates to zero,
whereas the product of their nonzero
saturated forms is $(0,\mathbf J)$.
There is nevertheless a precisely
defined unital connecting morphism.
For a finite-image label matrix let
$r_a(\Lambda)=\bigvee_b\Lambda_{ab}$,
$c_b(\Lambda)=\bigvee_a\Lambda_{ab}$,
and retain $\epsilon(\Lambda)=\bigvee_{a,b}\Lambda_{ab}$.
Then
\begin{equation}\tag{CFP46}
\begin{gathered}
 \mathfrak Z_H=\{(A,\Lambda):
       r_a(\Lambda)=c_b(\Lambda)=\epsilon(\Lambda)
                                  \text{ for every }a,b\},\\
 q_H:\mathfrak Z_H\twoheadrightarrow H\mathfrak M_L(\Gamma)H,
             \qquad M\longmapsto HMH,\\
 H(MN)H=(HMH)(HNH)\quad\Longleftrightarrow\quad
 \bigvee_b c_b(\Lambda)\wedge r_b(\Omega)
                  =\epsilon(\Lambda)\wedge\epsilon(\Omega),
 \qquad M=(A,\Lambda),\ N=(B,\Omega).
\end{gathered}
\end{equation}
All joins are again finite joins of
distinct palette values. The labels
of $MH$ and $HM$ at entry $(a,b)$
are respectively $r_a(\Lambda)$
and $c_b(\Lambda)$. They are equal
for every pair exactly when every
row and column join is the same
total join. Thus $\mathfrak Z_H$
is exactly the centralizer of $H$,
a unital subsemiring. For its
members one commutes the middle
$H$ past the factors and uses
$H^2=H$, proving multiplicativity
of $q_H$. It sends $I_{\mathrm{sp}}$
to the corner unit $H$, is onto,
and has the corner inclusion as
an additive and multiplicative
section, whose unit differs from
the ambient unit. For arbitrary
$M,N$, both compared amplitudes
are $AB$. Expanding the total
support join of their product
gives exactly the last displayed
criterion. This proves the full
morphism on its stated domain
and the precise obstruction to
extending its multiplicativity
to any additional pair. No
global-to-sparse relation has
been replaced by a claim of
nonidentity alone.

\begin{figure}[htbp]
\centering
$\begin{array}{ccc}
 \widehat\epsilon_q&\xrightarrow{\quad q\text{th power}\quad}&e\ne\tau\\
 {\scriptstyle\mathcal R_q}\downarrow&&\downarrow{\scriptstyle\mathcal R_q}\\
 e&\xrightarrow{\quad q\text{th power}\quad}&e
 \end{array}$
\caption{The finite characteristic-$p$ thickening with its
original support retained, CFP14--20. Reduction identifies
$\widehat\epsilon_q$ with $e$ by the exact equal-label
congruence; it does not set either element to $\tau$.
Its full prime-spectrum map is a homeomorphism despite
this nontrivial congruence. The source coordinates
$\epsilon_q=T_q-1$ and $\delta_q=1-T_q$ are related
with their sign in CFP1, CFP5 and CFP22, following
Connes--Consani--Marcolli, \emph{Fun with $\mathbb F_1$},
\texttt{filtration}, lines 1937--1953.}
\end{figure}
